\documentclass[11pt]{article}
\usepackage[margin=1in]{geometry}
\usepackage{amsmath,amssymb,mathtools}
\usepackage{booktabs,array}
\usepackage[hidelinks]{hyperref}
\newcommand{\NS}{\mathrm{NS}}
\newcommand{\R}{\mathrm{R}}
\newcommand{\Vir}{\mathrm{Vir}}
\newcommand{\SCA}{\mathrm{SCA}}
\newcommand{\cB}{\mathcal B}
\newcommand{\cR}{\mathcal R}
\newcommand{\eps}{\epsilon}
\newcommand{\ii}{\mathrm i}

\title{The Ramond $\Vir\oplus\Vir$ Branching Coefficient:\\
Direct Screening Matrix Elements and PBW Checks}
\author{Machine note}
\date{\today}

\begin{document}
\maketitle

\begin{abstract}
We formulate the Ramond analogue of the Hadasz--Jaskolski blow-up factor as
the matrix element of an explicit double-Virasoro primary in the free field.
The two Ramond spin fields turn its fermionic part into a zero-mode-resolved
Pfaffian.  For the first crossed state $(0,\frac34,\frac34)$ we evaluate this
matrix element directly on the two- and three-screening neutrality planes.
The two-screening Selberg integral produces the irreducible crossed
polynomial $H$, while the natural three-screening determinant produces the
factorized polynomial $K$.  We also retain a generic-momentum PBW/Ward path
sum as an independent finite-level oracle.  It is not itself a screening
derivation.  A generic-momentum, all-level Ramond screening formula remains
open.  The presently proved endpoint residues are rank-one maps on a
vector-valued contour problem; completion requires finite contour-orbit
closure, a full-rank set of collision and inversion/connection maps, base
vectors, and only then same-parity momentum interpolation.
\end{abstract}

\paragraph{Correction of scope.}
An earlier version called the PBW/Ward path sum below the final free-field
formula.  That was too strong: it is an exact basis reconstruction of the
three-form, hence an oracle and check, but not an evaluation of the Coulomb
matrix element.  The genuinely screened results of this version are stated
in Section~\ref{sec:direct-trial}.

\section{Conventions and statement of the problem}

We use the conventions of \texttt{Human Notes/SCblock.tex}.  Put
\begin{equation}
 Q=b+b^{-1},\qquad
 h_1=\frac18Q^2-\frac12P_1^2,\qquad
 h_j=\frac1{16}+\frac18Q^2-\frac12P_j^2\quad(j=2,3).
 \label{eq:weights}
\end{equation}
The ordered trinion is $(\NS,\R,\R)$, with the NS state in slot 1.  Its
branch labels obey
\begin{equation}
 n_1\in\tfrac12\mathbb Z,\qquad
 n_2,n_3\in\tfrac12\mathbb Z+\tfrac14.
 \label{eq:labels}
\end{equation}
It is useful to put $j_i=2n_i$.  Thus $j_1$ is integral and $j_2,j_3$ are
half-integral.  A Ramond summand occurs twice; we denote its copy by
$\eps_i\in\{0,1\}$.  The two chiral $(\NS,\R,\R)$ forms are
$\eta=\pm1$, and $f\in\{0,1\}$ denotes the parity of the complete
three-form.  In the human-note ground basis,
\begin{equation}
 \rho_0^{(\eta)}=\rho^{++}+\eta\rho^{--},\qquad
 \rho_1^{(\eta)}=\rho^{+-}+\ii\eta\rho^{-+}.
 \label{eq:human-forms}
\end{equation}
The normalized coefficient is
\begin{equation}
 \mathbb B_f^{(\eta)}
 =\frac{\widehat\rho_f^{(\eta)}
 (v_{1,n_1},v^{\eps_2}_{2,n_2},v^{\eps_3}_{3,n_3})}
 {\|v_{1,n_1}\|\,\|v^{\eps_2}_{2,n_2}\|\,
  \|v^{\eps_3}_{3,n_3}\|}.
 \label{eq:Bdef}
\end{equation}

\paragraph{Claim (the matrix element and its Selberg realization).}
Write $\boldsymbol n=(n_1,n_2,n_3)$ and
$\boldsymbol\eps=(\eps_2,\eps_3)$.  The quantity which the free-field
calculation has to produce is the raw double-Virasoro-primary matrix
element
\begin{equation}
 \mathsf M_{\boldsymbol\eps;f}^{(\eta)}(\boldsymbol P;\boldsymbol n)
 :=\left\langle v_{1,n_1}(P_1)\right|
 V_f^{(\eta)}\!\left[v_{2,n_2}^{\eps_2}(P_2)\right](1)
 \left|v_{3,n_3}^{\eps_3}(P_3)\right\rangle
 =\widehat\rho_f^{(\eta)}
 (v_{1,n_1},v_{2,n_2}^{\eps_2},v_{3,n_3}^{\eps_3}).
 \label{eq:matrix-element-claim}
\end{equation}
This is not a matrix element of the bare exponential alone: the three
states in \eqref{eq:matrix-element-claim} are the explicit simultaneous
$\Vir\oplus\Vir$ primaries, including their consecutive $\chi=f-\ii\psi$
strings and, on the Ramond legs, the chosen zero-mode components.

Let $\mathsf g_{\eps_2\eps_3;f}^{(\eta)}$ be the raw ground-spin matrix
in that component.  On a charge-neutral free-field representative with
$N$ insertions of the $b$-screening current, the claim is
\begin{equation}
 \left.\mathsf M_{\boldsymbol\eps;f}^{(\eta)}\right|_N
 =\mathsf g_{\eps_2\eps_3;f}^{(\eta)}
 \frac{\mathfrak S_N[
 F^{(\eta)}_{\boldsymbol n,\boldsymbol\eps};
 -b\alpha_1-u,-b\alpha_2-v,g]}
 {\mathfrak S_N[
 F_0^{(\eta,f)};
 -b\alpha_1-u_0,-b\alpha_2-v_0,g]},
 \qquad g=-\frac{bQ}{2}.
 \label{eq:matrix-element-Selberg-claim}
\end{equation}
Here $\mathfrak S_N$ is the Selberg functional
\eqref{eq:Selberg-functional}; $F_{\boldsymbol n,\boldsymbol\eps}$ is the
symmetric polynomial obtained by inserting the three branch-state strings
in the two-spin Pfaffian, while $F_0$ is obtained by removing those strings
but keeping the same Ramond ground component.  Equivalently, after the
ground-spin matrix is stripped off, the descendant-to-primary matrix-element
ratio is precisely the ratio of the two Selberg integrals in
\eqref{eq:matrix-element-Selberg-claim}.  Common ordered-contour phases and
factorials cancel in this ratio.

Combining \eqref{eq:Bdef} and
\eqref{eq:matrix-element-Selberg-claim} gives, on that neutrality plane,
\begin{equation}
 \left.\mathbb B_f^{(\eta)}\right|_N
 =\frac{\mathsf g_{\eps_2\eps_3;f}^{(\eta)}}
 {\|v_{1,n_1}\|\,\|v_{2,n_2}^{\eps_2}\|\,
  \|v_{3,n_3}^{\eps_3}\|}
 \frac{\mathfrak S_N[F_{\boldsymbol n,\boldsymbol\eps}^{(\eta)};
 -b\alpha_1-u,-b\alpha_2-v,g]}
 {\mathfrak S_N[F_0^{(\eta,f)};
 -b\alpha_1-u_0,-b\alpha_2-v_0,g]}.
 \label{eq:branching-Selberg-claim}
\end{equation}
The pole-cleared quantity used later is merely
$\cR|_N=(\prod_i\delta_{n_i}(P_i))\,\mathsf M|_N$; the $\delta$ factors
are not part of the physical three-point matrix element.  Also, $N$ here
is the number of screening currents, hence the dimension of the integral,
not a branch level.  A vector of contour-orbit integrals and its projectors
may be useful for evaluating the scalar ratio in
\eqref{eq:matrix-element-Selberg-claim}, but that vector is computational
auxiliary data rather than an additional branching coefficient.  An
efficient evaluation algorithm is deferred to the later contour-recursion
analysis.

Consequently a function of only $(b,P_i,n_i)$ is not defined in the R
sector: one must also say which $(\eps_2,\eps_3;f,\eta)$ component is meant.
Moreover, an unsquared coefficient changes under independent sign choices
for the branch highest vectors.  We therefore report the root-independent
quantity $(\mathbb B_f^{(\eta)})^2$; an unsquared representative is
fixed by the charge-basis convention below.

One Ramond summand starts at relative conformal weight
$2n^2-1/16$.  The exponent $2n^2-1/8$ in the theta-channel sewing formula
also subtracts the auxiliary Ramond vacuum weight.  These two numbers should
not be confused in a level count.

\section{The \texorpdfstring{$\ell$}{ell} products and the NS result}

For an integer $m>0$ define
\begin{equation}
 \ell(x,m)=
 \begin{cases}
 \displaystyle\prod_{\substack{r,s\geq0,\ r+s<m\\r+s\equiv0\ (2)}}
 (x+rb+sb^{-1}),&m\text{ even},\\[8pt]
 \displaystyle2^{1/8}\prod_{\substack{r,s\geq0,\ r+s<m\\r+s\equiv1\ (2)}}
 (x+rb+sb^{-1}),&m\text{ odd},
 \end{cases}
 \label{eq:ell-positive}
\end{equation}
and
\begin{equation}
 \ell(x,0)=1,\qquad
 \ell(x,m)=
 \begin{cases}
 (-1)^{-m/2}\ell(Q-x,-m),&m<0\text{ even},\\
 \ell(Q-x,-m),&m<0\text{ odd}.
 \end{cases}
 \label{eq:ell-negative}
\end{equation}
For later use set
\begin{equation}
 D_n(P)=\ell(2P,4n)\ell(Q+2P,4n),\qquad
 \delta_n(P)=2^{-\mathbf1_{4n\ {\rm odd}}/8}\ell(Q+2P,4n).
 \label{eq:leg-factors}
\end{equation}

We first reconstruct the part of the Hadasz--Jaskolski proof~\cite{HJ} that
has to be modified.  Let $f$ be the auxiliary Majorana fermion and $\psi$ the
super-Liouville fermion.  With their normalization,
\begin{equation}
 \chi=f-\ii\psi,\qquad
 |P,j\rangle=\chi_{-(2j-1)/2}\cdots\chi_{-1/2}|P,0\rangle
 \quad(j\in\mathbb Z_{>0}).
 \label{eq:NS-sea}
\end{equation}
Multiplication by the pole-clearing factor $\Omega(P,j)$ makes every PBW
coefficient polynomial in momentum.  A degree estimate gives degree at most
$j_1^2+j_2^2+j_3^2$ (one less for the odd three-form).

The two commuting screening charges are
\begin{equation}
 Q_b=\oint dz\,\psi(z)E^b(z),\qquad
 Q_{1/b}=\oint dz\,\psi(z)E^{1/b}(z).
 \label{eq:screenings}
\end{equation}
Reflection of any of the three Fock representations gives eight free-field
representatives of the same polynomial.  In the unreflected NS vacuum,
\begin{equation}
 \langle\chi(z)\chi(w)\rangle=0,
 \qquad \langle\chi(z)\psi(w)\rangle=-\frac{\ii}{z-w}.
 \label{eq:NS-zero}
\end{equation}
If there are fewer screenings than sea fermions, the numerator therefore
vanishes.  The eight representatives generate four complete families of
linear zeros; the degree bound shows that no momentum-dependent factor is
left over.

To determine the remaining constant, choose $N=j_1+j_2+j_3$ $b$-screenings.
Every sea fermion must end on a screening, and the fermionic numerator is
uniquely
\begin{equation}
 (-\ii)^N\prod_{a=1}^Nt_a^{-j_1}(1-t_a)^{-j_2}
 \prod_{a<c}(t_a-t_c).
 \label{eq:NS-vandermonde}
\end{equation}
The numerator is now an ordinary Selberg integral.  The denominator is the
Pfaffian Selberg integral evaluated in Appendix B of the paper.  Their ratio
fixes
\begin{equation}
 C^{j_3}_{j_2j_1}(b)=
 (-1)^{(j_1+j_2-j_3)/2}2^{(j_1+j_2+j_3)/2}.
 \label{eq:HJconstant}
\end{equation}
Translating to the human-note momentum and normalization convention gives,
for positive $n_i$,
\begin{align}
 B_a^{\NS}(P_i;n_i)
 &=\frac{(-1)^{2n_3}(-\ii)^a}{\prod_{i=1}^3\sqrt{D_{n_i}(P_i)}}
 \ell\left(\tfrac Q2+P_1+P_2+P_3,2(n_1+n_2+n_3)\right)\notag\\
 &\quad\times
 \ell\left(\tfrac Q2-P_1+P_2+P_3,2(-n_1+n_2+n_3)\right)\notag\\
 &\quad\times
 \ell\left(\tfrac Q2+P_1+P_2-P_3,2(n_1+n_2-n_3)\right)\notag\\
 &\quad\times
 \ell\left(\tfrac Q2-P_1+P_2-P_3,2(-n_1+n_2-n_3)\right).
 \label{eq:NS-result}
\end{align}
Negative labels are defined by $v_n(P)=v_{-n}(-P)$.  This derivation, rather
than merely the final product, is what we now generalize.

\section{Ramond branch vectors in the free field}

The R-sector decomposition proved by Bershtein--Shchechkin~\cite{BS} supplies the
missing highest vectors.  For $n>0$ put
\begin{equation}
 M(n)=2n-\frac12\in\mathbb Z_{\geq0}.
 \label{eq:Mdef}
\end{equation}
The following literal recurrence fixes the two copies and all zero-mode
phases.  It is the convention used by the code.  Start with the tensor
ground $|0,0\rangle$ and the ordered list
\begin{equation}
 \mathbf o_n=(0,-1,\ldots,-M).
 \label{eq:R-operator-list}
\end{equation}
If $\eps\ne M+1\pmod2$, append a tagged zero $0^\star$ to this list.  Read
the resulting list from right to left.  At each entry $r$, replace every
current tensor state by the two terms
\begin{equation}
 |A\rangle_F\otimes|X\rangle_R\longmapsto
 f_r|A\rangle_F\otimes|X\rangle_R
 -\ii(-1)^{|A|}|A\rangle_F\otimes\psi_r|X\rangle_R.
 \label{eq:R-binary-update}
\end{equation}
For the tagged entry only, use $\psi_{0^\star}=-\psi_0$ in the second
term.  The Clifford action is
$f_0|a\rangle=2^{-1/2}|1-a\rangle$ and
$\psi_0|g\rangle=2^{-1/2}|1-g\rangle$, with the usual signs from moving a
zero mode through existing creators.  The final finite sum is, by
definition, $W_n^{(\eps)}$.  For example,
\begin{align}
 W_{1/4}^{(0)}&=|0,0\rangle+\ii|1,1\rangle,\notag\\
 W_{1/4}^{(1)}&=2^{-1/2}(|1,0\rangle-\ii|0,1\rangle).
 \label{eq:R-ground-branches}
\end{align}
Thus the one-zero-mode string is the $\eps=1$ copy at $n=1/4$; the
$\eps=0$ copy contains the tagged opposite zero mode.  This corrects the
potentially ambiguous shorthand ``one copy plus the other zero mode.''
The states produced by \eqref{eq:R-binary-update} are simultaneous highest
vectors of the two Virasoro algebras.

There is an exact finite definition of their norms that does not rely on an
all-level product conjecture.  Expand the finite string in an auxiliary
fermion basis and a Ramond PBW basis,
\begin{equation}
 W_n^{(\eps)}=\sum_{(A,B,\mu)\in\mathcal I_n^{(\eps)}}
 C_{A,B,\mu}^{(\eps)}(P,Q)\,|A\rangle_F\otimes|B,\mu\rangle_R.
 \label{eq:R-PBW-expansion}
\end{equation}
The coefficients are fixed uniquely by the triangular free-field formulas
for $L_m,G_r$ and the chosen $G_0$ action \eqref{eq:G0}.  If $G^F$ and $G^R$
are the finite auxiliary and Ramond Gram matrices at the required level, put
\begin{equation}
 \mathcal N_R^{(\eps)}(n,P)=
 \sum C_{A,B,\mu}^{(\eps)}C_{A',B',\nu}^{(\eps)}
 G^F_{AA'}G^R_{B\mu,B'\nu}.
 \label{eq:R-exact-norm}
\end{equation}
This is a finite rational formula for every fixed $n$.

The factorized all-level norm suggested by the reflection zeros and the
Ramond Whittaker coefficient is
\begin{align}
 \|W_n^{(0)}\|^2&=c_R^{(0)}(n)
 \frac{\ell(2P,4n)}{\ell(Q+2P,4n)},&
 c_R^{(0)}(n)&=2^{2\lfloor M/2\rfloor+1},\notag\\
 \|W_n^{(1)}\|^2&=c_R^{(1)}(n)
 \frac{\ell(2P,4n)}{\ell(Q+2P,4n)},&
 c_R^{(1)}(n)&=-2^{2\lfloor(M+1)/2\rfloor}.
 \label{eq:R-product-norm}
\end{align}
The same oscillator argument that proves the NS pole-clearing statement
gives the common leg denominator $\delta_n(P)$ in \eqref{eq:leg-factors}.
Equation \eqref{eq:R-product-norm} is verified directly for the states needed
below, $n=1/4,3/4,5/4$.  Its continuation is the natural all-level product
conjecture; it is closely related to the still-conjectural Ramond Whittaker
coefficient in Ref.~\cite{BS}.  The exact result below instead uses
\eqref{eq:R-exact-norm}; the compact $\ell$-only corollary uses
\eqref{eq:R-product-norm}.

For the bosonization paragraph only, use canonically normalized complex
fermions $\widehat\chi^\pm=(f\pm\ii\psi)/\sqrt2$.  Then
\begin{equation}
 \widehat\chi^\pm=\kappa_\pm:e^{\pm\ii H}:,\qquad
 \psi=\frac{\widehat\chi^+-\widehat\chi^-}{\ii\sqrt2}
 \label{eq:bosonization}
\end{equation}
which makes the
new feature manifest.  A screening current is the sum of two $H$-charges.
In the NS calculation \eqref{eq:NS-zero} forced a unique colour.  Between
two Ramond spin fields, charge may also run through the Ramond ground
doublet.  The screened numerator is consequently a sum over colours.  This
is the precise sense in which the R problem is matrix-valued.

\section{The finite Ramond screening kernel}
\label{sec:kernel}

We give a closed definition that involves only finite sums, Pfaffians and
gamma functions.  It is also an algorithm for arbitrary fixed branch
labels.

\subsection{Free-fermion part}

Move the two Ramond punctures to $0$ and $\infty$.  In either Majorana
factor the even Ramond propagator is
\begin{equation}
 K_R(z,w)=\frac{z+w}{2\sqrt{zw}(z-w)}.
 \label{eq:R-kernel}
\end{equation}
An even matrix element is the Pfaffian of $K_R(z_a,z_b)$.  An odd matrix
element is the Pfaffian of the augmented antisymmetric matrix whose extra
row is $2^{-1/2}z_a^{-1/2}$.  The signs of that row are fixed by
\begin{equation}
 G_0w^+=\ii\beta e^{-\ii\pi/4}w^-,\qquad
 G_0w^-=\ii\beta e^{\ii\pi/4}w^+.
 \label{eq:G0}
\end{equation}

Here is the mode-extraction convention, which removes another possible
ambiguity.  Put the punctures at $(\infty,1,0)$ and choose local coordinates
\begin{equation}
 z_1(x)=x^{-1},\qquad z_2(x)=1+x,\qquad z_3(x)=x.
 \label{eq:local-coordinates}
\end{equation}
For a weight-$1/2$ field $X$, let
$X^{[i]}(x)=(dz_i/dx)^{1/2}X(z_i(x))$, with square-root lifts
$(\varsigma_1,\varsigma_2,\varsigma_3)=(-1,\ii,1)$.  A creator $X_{-m}$ on
leg $i$ means the literal coefficient extraction
\begin{equation}
 [X_{-m}]_i\Phi=\mathop{\rm Res}_{x=0}
 x^{-m-1/2}\Phi(\ldots,X^{[i]}(x),\ldots).
 \label{eq:mode-extraction}
\end{equation}
At infinity the string is BPZ reversed and $X_r^\dagger=-X_{-r}$.
The auxiliary and physical ground matrices are, respectively,
\begin{equation}
 \Gamma_F^{0}=\begin{pmatrix}1&0\\0&-1\end{pmatrix},\quad
 \Gamma_F^{1}=\begin{pmatrix}0&1\\-1&0\end{pmatrix},\qquad
 \Gamma_\psi^{\eta,0}=\begin{pmatrix}1&0\\0&\eta\end{pmatrix},\quad
 \Gamma_\psi^{\eta,1}=\begin{pmatrix}0&1\\\ii\eta&0\end{pmatrix}.
 \label{eq:ground-matrices}
\end{equation}
For any ordered list of Majorana fields, write
$\mathcal W_\Gamma$ for Wick contraction using $K_R$, the augmented row
above when the list is odd, and the relevant entry of $\Gamma$.  Thus no
three-point constant is left inside $\mathcal W_\Gamma$.  For the master
component $(\eps_3,f)=(0,0)$ use
$\Gamma_F^{a_F}$ with $a_F=2n_1+\eps_2\pmod2$ and
$\Gamma_\psi^{\eta,0}$.

Expand every ordered string \eqref{eq:NS-sea}, \eqref{eq:R-binary-update} by choosing,
for each $\chi=f-\ii\psi$, its auxiliary or physical colour.  Let
$a_i$ denote such a finite binary path on leg $i$, $C_i(a_i)$ its sign and
power of $\ii$, and $p_F(a_i),p_\psi(a_i)$ its two parities.  If
$\mathbf f(a)$ and $\boldsymbol\psi(a)$ are the auxiliary and physical
external mode lists, including their ground indices, define the Laurent
polynomial by
\begin{equation}
 \mathsf P^{(\eta)}_{a_1a_2a_3}(\boldsymbol t)
 =\operatorname{ME}_{a_1a_2a_3}\!\left[
 \mathcal W_{\Gamma_F}(\mathbf f(a))\,
 \mathcal W_{\Gamma_\psi^{\eta}}
 (\boldsymbol\psi(a),\psi(t_1),\ldots,\psi(t_N))\right],
 \label{eq:P-explicit}
\end{equation}
where $\operatorname{ME}$ is the ordered product of residues
\eqref{eq:mode-extraction}.  The full fermionic
numerator is
\begin{equation}
 \mathsf H^{(\eta)}_{\boldsymbol n,\boldsymbol\eps}(\boldsymbol t)
 =\sum_{a_1,a_2,a_3}
 (-1)^{p_\psi(a_1)[p_F(a_2)+p_F(a_3)]+p_\psi(a_2)p_F(a_3)}
 \prod_{i=1}^3 C_i(a_i)\,
 \mathsf P^{(\eta)}_{a_1a_2a_3}(\boldsymbol t).
 \label{eq:fermion-path-sum}
\end{equation}
There are no fitted coefficients in \eqref{eq:fermion-path-sum}; its number
of summands is finite, at most $2^{j_1+M(n_2)+M(n_3)+2}$ before zero terms
are removed.  The displayed exponent is increased by one for each tagged
opposite zero mode.  The $\eta$ projection is exactly
\eqref{eq:human-forms}.

After pulling out the endpoint powers required by the OPE, the only
coincident-point singularity is the simple free-fermion pole.  Hence
antisymmetry in the screening coordinates implies
\begin{equation}
 \mathsf H^{(\eta)}(\boldsymbol t)
 =\prod_a t_a^{-u}(1-t_a)^{-v}\,\Delta(\boldsymbol t)^{-1}\,
 F^{(\eta)}_{\boldsymbol n,\boldsymbol\eps}(\boldsymbol t),
 \qquad \Delta(\boldsymbol t)=\prod_{a<c}(t_a-t_c),
 \label{eq:symmetric-screening-polynomial}
\end{equation}
where $F^{(\eta)}$ is a symmetric polynomial.  In the one-colour NS
numerator \eqref{eq:NS-vandermonde}, $F=\Delta^2$, so the effective
Vandermonde exponent is shifted upward by two.  Direct contractions between
the two Ramond strings supply the other terms of the Ramond $F$.

There is no implicit choice in the endpoint exponents in
\eqref{eq:symmetric-screening-polynomial}.  Given the explicitly extracted
Laurent polynomial $\mathsf H$, define $u$ to be the negative of its least
power in any $t_a$.  After removing $\prod_at_a^{-u}$ and replacing
$t_a=1-s_a$, define $v$ in the same way at $s_a=0$.  Fermionic
antisymmetry and the simple-pole bound then make
$F=\Delta\prod_at_a^u(1-t_a)^v\mathsf H$ a polynomial.  This finite
valuation prescription can therefore be implemented without making a
factorization ansatz for $F$.

\subsection{Evaluation by a finite Jack--Selberg sum}

For a symmetric polynomial $F$ define
\begin{equation}
 \mathfrak S_N[F;A,B,g]
 =\int_{[0,1]^N}\prod_{a=1}^Ndt_a\,
 t_a^A(1-t_a)^B|\Delta(\boldsymbol t)|^{2g}F(\boldsymbol t).
 \label{eq:Selberg-functional}
\end{equation}
Expand, finitely,
\begin{equation}
 F=\sum_\lambda f_\lambda P_\lambda^{(1/g)},\qquad
 [x]_{\lambda}^{(g)}=
 \prod_{r=1}^N\bigl(x+(1-r)g\bigr)_{\lambda_r}.
 \label{eq:Jack-expansion}
\end{equation}
Kadell's Selberg average evaluates \eqref{eq:Selberg-functional} as
\begin{equation}
 \mathfrak S_N[F;A,B,g]
 =S_N(A,B;g)\sum_\lambda f_\lambda
 P_\lambda^{(1/g)}(1^N)
 \frac{[A+1+(N-1)g]_{\lambda}^{(g)}}
 {[A+B+2+2(N-1)g]_{\lambda}^{(g)}},
 \label{eq:Jack-Selberg}
\end{equation}
where
\begin{equation}
 S_N(A,B;g)=\prod_{r=0}^{N-1}
 \frac{\Gamma(A+1+rg)\Gamma(B+1+rg)
 \Gamma(1+(r+1)g)}
 {\Gamma(A+B+2+(N+r-1)g)\Gamma(1+g)}.
 \label{eq:Selberg}
\end{equation}
Thus \eqref{eq:Jack-Selberg} is a finite sum of explicit gamma and
Pochhammer products.  It is not an unevaluated contour integral.  As in the
original Coulomb-gas calculation, the integral is first defined in its
convergence chamber and the displayed gamma-function expression gives its
meromorphic continuation to the physical value of $g$.

At a pure $b$-screening neutrality point
\begin{equation}
 \alpha_3=\alpha_1+\alpha_2+Nb,\qquad
 \alpha_i=\frac Q2+P_i,
 \label{eq:neutrality}
\end{equation}
the bosonic contractions give $A=-b\alpha_1-u$ and
$B=-b\alpha_2-v$.  A second, dual evaluation uses
$\alpha_3=\alpha_1+\alpha_2+N/b$ and $b\mapsto b^{-1}$.  For the first
chamber, the common Pfaffian-Selberg exponent is
\begin{equation}
 g=-\frac{bQ}{2}=-\frac{b^2+1}{2}.
 \label{eq:g}
\end{equation}
For the one-colour insertion $F=\Delta^2$, this becomes the ordinary
Selberg exponent $g+1=(1-b^2)/2$, reproducing the numerator integral used
by Hadasz--Jaskolski.
The numerator of the normalized screened matrix element is therefore
\eqref{eq:Jack-Selberg} with the polynomial obtained from
\eqref{eq:fermion-path-sum}.  Its denominator is the same expression with
all three branch strings removed but with the same R ground component.
More explicitly, let $\mathsf g_{\eps_2\eps_3;f}^{(\eta)}$ be the raw
charge-basis ground amplitude for the selected component.  If
$F_0^{(\eta,f)}$ and $(u_0,v_0)$ denote the corresponding ground polynomial
and its endpoint valuations, then at the neutrality point
\eqref{eq:neutrality} we define
\begin{equation}
 \cR_{\eps_2\eps_3;f}^{(\eta)}\big|_N
 =\mathsf g_{\eps_2\eps_3;f}^{(\eta)}
 \delta_{n_1}(P_1)\delta_{n_2}(P_2)\delta_{n_3}(P_3)
 \frac{\mathfrak S_N[F^{(\eta)}_{\boldsymbol n,\boldsymbol\eps};
 -b\alpha_1-u,-b\alpha_2-v,g]}
 {\mathfrak S_N[F_0^{(\eta,f)};
 -b\alpha_1-u_0,-b\alpha_2-v_0,g]}.
 \label{eq:Rpoly-neutral}
\end{equation}
The ordered-contour factorial is common to numerator and denominator and
cancels in this ratio.  The prefactor restores the ground normalization
removed by forming the ratio; in particular
$\mathsf g_{00;0}^{(\eta)}=1+\ii\eta$ and
$\mathsf g_{10;0}^{(\eta)}=-e^{-\ii\eta\pi/4}$ in the convention
\eqref{eq:G0}.  We use the shorter notation
\begin{equation}
 \cR_{\eps_2\eps_3;f}^{(\eta)}
 (P_1,P_2,P_3;n_1,n_2,n_3).
 \label{eq:Rpoly-def}
\end{equation}
Equations \eqref{eq:fermion-path-sum}--\eqref{eq:Rpoly-neutral} give a
closed finite evaluation of the native ground-resolved Coulomb functional.
It is a fixed SCA chiral form only when the screening holonomy selects that
form.  Applying a constant ground-space rotation to the other excited
component is not justified; the reflected vertex must be evaluated
independently.

For completeness, a conditional screening reconstruction away from a
neutrality hyperplane follows if the pole-cleared matrix element is a
polynomial with the following sharp degree.  For the even component put
\begin{equation}
 d=j_1^2+j_2^2+j_3^2-\frac12,
 \label{eq:degree}
\end{equation}
and the odd component has degree at most $d-1$.  Define $N_0$ to be the
smallest element of $\{0,1\}$ for which the ground denominator in
\eqref{eq:Rpoly-neutral} is nonzero for the selected component.  Choose
$d+1$ neutrality points of that parity,
$N=N_0,N_0+2,\ldots,N_0+2d$.  If their third momenta are
$P_3^{(a)}$, then
\begin{equation}
 \cR(P_3)=\sum_{a=0}^{d}\cR(P_3^{(a)})
 \prod_{\substack{c=0\\c\ne a}}^d
 \frac{P_3-P_3^{(c)}}{P_3^{(a)}-P_3^{(c)}}.
 \label{eq:Lagrange}
\end{equation}
If every node in \eqref{eq:Lagrange} is independently evaluated as the same
analytic SCA form, this would define it at generic
$(P_1,P_2,P_3)$.  That premise is currently missing for the complementary
Ramond charge chart: the reflected fermion contains bosonic-current modes.
The degree bound is supported by clearing the reflection denominators in
the finite PBW oracle---leg $i$ contributes $j_i^2$, while the two Ramond
zero-mode seas remove one half in total---but a degree bound does not
supply the absent screening nodes.

The logical status should be stated precisely.  Polynomiality and this
degree argument are the all-level step needed to promote the screened
values \eqref{eq:Rpoly-neutral} to the generic value \eqref{eq:Lagrange}.
They are proved in the NS case and verified for the R states used in
Section~\ref{sec:low}; a general R proof is not supplied by the cited
literature.  Hence the screening-only interpolation is a concrete
all-level conjecture, not by itself the unconditional definition below.

\subsection{Generic-momentum PBW/Ward oracle}

For comparison we now give a finite evaluator with no screening number or
interpolation nodes.  This is the object against which the contour
calculation is tested; it is not used to derive a screened value.  At each
physical level let $T_\sigma(P,Q)$ be the
square matrix whose columns are SCA PBW words expanded in the free Fock
basis using, for negative modes,
\begin{align}
 L_k^{(\sigma)}&=\frac12\sum_{m\ne0,k}:c_{k-m}c_m:
 +\frac12\sum_r r:\psi_{k-r}\psi_r:
 +\frac{\ii}{2}(kQ+2\sigma P)c_k,\notag\\
 G_r^{(\sigma)}&=\sum_{m\ne0}c_m\psi_{r-m}
 +\ii(rQ+\sigma P)\psi_r.
 \label{eq:free-field-transition}
\end{align}
For positive branch labels $\sigma=-1$; reflection gives the other sheet.
Expand the NS sea \eqref{eq:NS-sea} and the R recurrence
\eqref{eq:R-binary-update}, and multiply every physical Fock endpoint by
$T_{-}^{-1}$.  Denote the resulting finite component lists by
$\mathcal C_i$.  An element $x_i\in\mathcal C_i$ consists of an auxiliary
Fock state $A_i$, a physical PBW state $V_i$, and its rational coefficient
$c_i(P_i,Q)$.

Let $\rho_F^{a_F}$ be the auxiliary Majorana three-form fixed by
\eqref{eq:R-kernel}, \eqref{eq:mode-extraction}, and
$\Gamma_F^{a_F}$.  Let $\rho_{\SCA}^{\eta,0}$ be the unique NS--R--R SCA
three-form obtained from the SCA commutators and the ground matrix
$\Gamma_\psi^{\eta,0}$.  Put
\begin{gather}
 a_F=2n_1+\eps_2\pmod2,\qquad
 \eta_{\rm eff}=(-1)^{2n_1}\eta,\notag\\
 \mathcal A_{\eps_2}^{(\eta)}(b,\boldsymbol P;\boldsymbol n)
 =\sum_{x_i\in\mathcal C_i}(-1)^{\kappa(x_1,x_2,x_3)}
 c_1c_2c_3\,
 \rho_F^{a_F}(A_1,A_2,A_3)\,
 \rho_{\SCA}^{\eta_{\rm eff},0}(V_1,V_2,V_3),
 \label{eq:literal-A}\\
 \kappa=p_\psi(x_1)[p_F(x_2)+p_F(x_3)]
 +p_\psi(x_2)p_F(x_3)\pmod2.
 \notag
\end{gather}
All three sums in \eqref{eq:literal-A} are finite, so this is an unambiguous
PBW/Ward evaluator for every fixed set of labels; it is implemented in
\path{Code/screening_kernel.py}.  It is nevertheless tautological from the
viewpoint of the desired free-field calculation, because
$\rho_{\SCA}$ is left to the Ward recursion.  The pole-cleared notation used
above is only the abbreviation
\begin{equation}
 \cR_{\eps_2 0;0}^{(\eta)}
 =\delta_{n_1}(P_1)\delta_{n_2}(P_2)\delta_{n_3}(P_3)
 \mathcal A_{\eps_2}^{(\eta)}.
 \label{eq:R-is-A}
\end{equation}
Equation \eqref{eq:Rpoly-neutral} is the proposed screening-charge
evaluation of this already-defined function on a neutrality hyperplane.

\subsection{Normalized branching coefficient}

It is convenient to present the master component with $\eps_3=f=0$ first.
Let
\begin{equation}
 c_{\NS}(n)=(-1)^{2n}2^{2n},\qquad
 \mathcal N_{\NS}(n,P)=c_{\NS}(n)
 \frac{\ell(2P,4n)}{\ell(Q+2P,4n)},
 \label{eq:cNS}
\end{equation}
and use the exact finite R norm \eqref{eq:R-exact-norm}.  In the raw
charge-basis normalization of \eqref{eq:NS-sea} and
\eqref{eq:R-binary-update}, the PBW/Ward oracle for the matrix element is
\begin{equation}
 \widehat\rho_0^{(\eta)}
 =\mathcal A_{\eps_2}^{(\eta)}
 =\frac{\cR_{\eps_2 0;0}^{(\eta)}}
 {\delta_{n_1}(P_1)\delta_{n_2}(P_2)\delta_{n_3}(P_3)}.
 \label{eq:raw-rho}
\end{equation}
Combining \eqref{eq:raw-rho} with the three raw norms gives the exact
normalization identity
\begin{equation}
 \boxed{
 \left(\mathbb B_f^{(\eta)}
 [\boldsymbol n;\eps_2,\eps_3]\right)^2
 =(-1)^{\eps_3}(-\ii)^f
 \frac{\left(\mathcal A_{\eps_2}^{(\eta)}\right)^2}
 {\mathcal N_{\NS}(n_1,P_1)
  \mathcal N_R^{(\eps_2)}(n_2,P_2)
  \mathcal N_R^{(0)}(n_3,P_3)}.}
 \label{eq:master-B-exact}
\end{equation}
Here $\mathcal A$ is the PBW/Ward sum \eqref{eq:literal-A}.  Equation
\eqref{eq:master-B-exact} becomes a free-field result only after the same
raw matrix element has been obtained from the screening contour.

If the factorized norm conjecture \eqref{eq:R-product-norm} is used,
\eqref{eq:master-B-exact} reduces to the compact $\ell$-only expression
\begin{equation}
 \boxed{
 \left(\mathbb B_f^{(\eta)}
 [\boldsymbol n;\eps_2,\eps_3]\right)^2
 =(-1)^{\eps_3}(-\ii)^f\,
 \frac{2^{1/2}\left(\cR_{\eps_2 0;0}^{(\eta)}\right)^2}
 {c_{\NS}(n_1)c_R^{(\eps_2)}(n_2)c_R^{(0)}(n_3)
  \prod_{i=1}^3D_{n_i}(P_i)}.}
 \label{eq:master-B}
\end{equation}
Under both conjectural replacements just stated---the factorized R norm and
the generic-momentum screening interpolation---all the continuous dependence
on $b,P_i,n_i$ in \eqref{eq:master-B} lies in the finite polynomial
$\cR$ defined by \eqref{eq:fermion-path-sum}--\eqref{eq:Lagrange}.  The
$\eps_3$ and $f$ dependence in \eqref{eq:master-B-exact} is universal.  The
$\eps_2$ dependence is retained inside the screening polynomial; it is not
legitimate to discard it before doing the zero-mode projection.
Negative branch labels require no new formula.  Reflect one leg at a time,
\begin{equation}
 (n_i,P_i)\longmapsto(-n_i,-P_i),
 \label{eq:R-reflection}
\end{equation}
use \eqref{eq:ell-negative}, and evaluate the positive representative.
The possible sign of the reflected highest vector disappears from the
squared coefficient \eqref{eq:master-B-exact}.

\section{Relation to the four-\texorpdfstring{$\ell$}{ell} conjecture}

Define the scalar Hadasz--Jaskolski continuation
\begin{align}
 \mathcal P(\boldsymbol n;\boldsymbol P)
 &=\ell(\tfrac Q2+P_1+P_2+P_3,2(n_1+n_2+n_3))\notag\\
 &\quad\times\ell(\tfrac Q2-P_1+P_2+P_3,2(-n_1+n_2+n_3))\notag\\
 &\quad\times\ell(\tfrac Q2+P_1+P_2-P_3,2(n_1+n_2-n_3))\notag\\
 &\quad\times\ell(\tfrac Q2-P_1+P_2-P_3,2(-n_1+n_2-n_3)).
 \label{eq:Pcal}
\end{align}
The quantity that should be divided out, rather than subtracted, is
\begin{equation}
 U_{\eps_2}^{(\eta)}
 =\frac{\cR_{\eps_2 0;0}^{(\eta)}}
 {\cR_{\eps_2 0;0}^{(\eta)}\big|_{F=\Delta^2}}
 =\frac{\cR_{\eps_2 0;0}^{(\eta)}}
 {\hbox{the appropriate reflected-sheet version of }\mathcal P}.
 \label{eq:Ufactor}
\end{equation}
Here $F=\Delta^2$ means retaining only the path in which every external sea
fermion ends on a screening, as in \eqref{eq:NS-vandermonde}.  In the NS sector $U=1$
identically.  In the R sector it remains one whenever the charge projection
forces a single screening colour.  In general $U$ is a momentum-dependent
finite Jack sum.  It is not the field normalization $\Omega_k(\alpha)$ of
the double-Liouville dictionary~\cite{SS}, and it is not a
momentum-independent normalization constant.

\section{Low-level channels and direct screening restrictions}
\label{sec:low}

This section is placed after the derivation deliberately.

\subsection{Two Ramond ground branches}

For $(n_1,n_2,n_3)=(0,1/4,1/4)$ the symmetric screening polynomial is
constant.  In the raw convention,
\begin{equation}
 \cR_{00;0}^{(\eta)}=1+\ii\eta,\qquad
 \cR_{10;0}^{(\eta)}=-e^{-\ii\eta\pi/4}.
 \label{eq:ground-R}
\end{equation}
Using $c_R^{(0)}(1/4)=2$ and $c_R^{(1)}(1/4)=-1$ in
\eqref{eq:master-B} gives, for both $\eps_2$,
\begin{equation}
 \left(\mathbb B_f^{(\eta)}[0,\tfrac14,\tfrac14;
 \eps_2,\eps_3]\right)^2
 =(-1)^{\eps_3}(-\ii)^f\frac{\ii\eta}{2}.
 \label{eq:ground-B}
\end{equation}

\subsection{First crossed kernel}

The first nonconstant $F$ occurs when both Ramond strings contain a
nonzero mode, at $(0,3/4,3/4)$.  Put
\begin{gather}
 x_{++}=\frac Q2+P_1+P_2+P_3,\qquad
 x_{--}=\frac Q2+P_1-P_2-P_3,\\
 x_2=\frac Q2-P_1+P_2+P_3,\qquad E_j=Q+2P_j,\\
 a_j=(2P_j)^2+Q(2P_j)+1,\qquad
 d_j=E_j^2+QE_j+1,\\
 K=(x_{++}^2+Qx_{++}+1)(x_2^2+Qx_2+1),\qquad
 L=x_{++}(x_{--}-Q),\\
 H=L^2+2L(E_2E_3+1)+d_2d_3.
 \label{eq:KH}
\end{gather}
The one-colour component gives $K$.  The remaining contraction paths in
$F^{(-)}$ give the crossed result $H$, equivalently
\begin{equation}
 H=(1,L)
 \begin{pmatrix}d_2d_3&1+E_2E_3\\1+E_2E_3&1\end{pmatrix}
 \binom{1}{L}.
 \label{eq:hard-kernel}
\end{equation}
For the $\eps_2=0$ master the normalized answer is
\begin{equation}
 \left(\mathbb B_f^{(\eta)}
 [0,\tfrac34,\tfrac34;0,\eps_3]\right)^2
 =(-1)^{\eps_3}(-\ii)^f\frac{\ii\eta}{2}
 \frac{F_\eta^2}{a_2a_3d_2d_3},
 \qquad F_+=K,\quad F_-=H.
 \label{eq:hard-B}
\end{equation}
The $\eps_2=1$ component is obtained from the same literal sum
\eqref{eq:literal-A}; it is not implicitly identified with
\eqref{eq:hard-B}.
Since $H$ is irreducible over the rational function field in
$(Q,P_1,P_2,P_3)$, it cannot be a four-$\ell$ product on any reflected
momentum sheet.  This is not an obstruction to free fields; it is the first
place where the two-colour screening polynomial is nonconstant.  Relative
to the literal four-$\ell$ factor, $U_+=1$ and $U_-=H/K$.

\subsection{Direct screening matrix element on two hyperplanes}
\label{sec:direct-trial}

We now compute rather than reconstruct the first crossed three-form.  Keep
the puncture order $(\infty,1,0)=(\NS_1,\R_2,\R_3)$ and define
\begin{equation}
 \mathfrak M_f^{(\eta)}=
 \frac{
 \langle v_{1,0}(P_1)|
 V[v_{2,3/4}^{0}(P_2)](1)|v_{3,3/4}^{0}(P_3)\rangle_f^{(\eta)}
 }{
 \langle v_{1,0}(P_1)|V[w_2](1)|w_3\rangle_f^{(\eta)}}.
 \label{eq:direct-matrix-element}
\end{equation}
The numerator uses the literal consecutive $\chi=f-\ii\psi$ strings on
the two Ramond fields and $N$ insertions of
$\psi(t)E^b(t)$.  The denominator is the primary two-spin BFL integral in
the same ground channel.

At $N=2$, charge neutrality is
\begin{equation}
 P_1=-\frac Q2-P_2-P_3-2b.
 \label{eq:N2-plane}
\end{equation}
Writing $e_1=t_1+t_2$ and $e_2=t_1t_2$, direct mode extraction from the two
$\chi$ strings and the two-spin Pfaffian gives, for $f=0$ and $\eta=-1$,
\begin{equation}
 F_2(t_1,t_2)=\frac{1-\ii}{2}(2e_2-e_1)
 \bigl(e_2^2-e_1e_2+e_1^2-3e_2\bigr).
 \label{eq:N2-integrand}
\end{equation}
Expanding only in products of elementary symmetric functions and applying
their exact Selberg moments yields
\begin{align}
 \mathfrak M_0^{(-)}&=(-1+\ii)R_2,&
 \mathfrak M_1^{(-)}&=-\sqrt2\,\ii R_2,\notag\\
 R_2&=
 \frac{13b^4+10b^3(P_2+P_3)+4b^2P_2P_3+8b^2
       +4b(P_2+P_3)+4}
 {(b^2+2bP_2+2)(b^2+2bP_3+2)}.
 \label{eq:N2-result}
\end{align}
This computation imports no SCA three-form.  Symbolic reduction gives the
nontrivial identity
\begin{equation}
 R_2=\left.\frac{H}{d_2d_3}\right|_{\eqref{eq:N2-plane}}.
 \label{eq:N2-H}
\end{equation}
Thus the irreducible $H$ is directly visible in the screening integral, not
only in the PBW calculation.

At the natural $N=3$ node,
$P_1=-Q/2-P_2-P_3-3b$, the physical form is $\eta=+1$.  The external--screen
block reduces the Pfaffian to one determinant, followed by an ordinary
Selberg product.  The answer is
\begin{align}
 \mathfrak M_0^{(+)}&=-4(1+\ii)R_3,&
 \mathfrak M_1^{(+)}&=-4\sqrt2\,R_3,\notag\\
 R_3&=
 \frac{b^2(3b^2-1)
 (2b^2+bP_2+bP_3+1)(5b^2+2bP_2+2bP_3+1)}
 {(b^2+2bP_2+2)(b^2+2bP_3+2)
  (2b^2+2bP_2+1)(2b^2+2bP_3+1)}.
 \label{eq:N3-result}
\end{align}
Direct simplification gives
\begin{equation}
 4R_3=\left.\frac{K}{d_2d_3}
 \right|_{P_1=-Q/2-P_2-P_3-3b}.
 \label{eq:N3-K}
\end{equation}
Equations \eqref{eq:N2-H} and \eqref{eq:N3-K} are genuine contour checks
of the two channels.  They do not yet determine either generic degree-four
polynomial from screening data alone.  The executable derivation is
\path{Code/direct_matrix_element_trial.py}.

\section{Independent PBW/Ward comparison through onset level
\texorpdfstring{$7/2$}{7/2}}

The comparison uses explicit simultaneous $\Vir\oplus\Vir$ primaries and
the generalized $(\NS,\R,\R)$ Ward identities, not the screening contours.
The intrinsic NS-primary parity is mandatory: the Ward sign uses
\begin{equation}
 \eta_{\rm eff}=(-1)^{p_\phi}\eta,
 \qquad p_\phi=2n_1\pmod2.
 \label{eq:eta-effective}
\end{equation}
The labels whose individual branch onset does not exceed $7/2$ are
\begin{equation}
 n_1\in\{0,\tfrac12,1\},\qquad
 n_2,n_3\in\{\tfrac14,\tfrac34,\tfrac54\}.
 \label{eq:grid}
\end{equation}
There are $27$ level triples, four independent continuous masters
$(\eps_2,\eta)$ per triple, and sixteen discrete restrictions after
$(\eps_3,f)$ are restored.  The audit checks the following separately:
\begin{enumerate}
 \item the raw R norms at $n=1/4,3/4,5/4$;
 \item all four PBW/path master amplitudes at each of the 27 triples;
 \item the universal reduction
 $(\mathbb B_f^{(\eta)})^2=(-1)^{\eps_3}(-\ii)^f
 (\mathbb B_0^{(\eta)}[\eps_2,0])^2$;
 \item the scalar-product specialization and its first failure
 \eqref{eq:hard-B}.
\end{enumerate}
The exact arithmetic report and the commands used to regenerate it are in
\path{Code/verification_report.md}.  The continuous grid audit constructs
the branch strings in two different ways (finite colour paths versus direct
simultaneous-primary assembly); both use the same independently audited Ward
functional.  It is therefore a check of the literal finite function,
normalization and phase restoration, not a proof of the generic-momentum
screening interpolation or of the Jack--Selberg identity.  The comparison
confirms \eqref{eq:literal-A} throughout \eqref{eq:grid}.  The scalar
four-$\ell$ specialization matches only the one-colour components; the
crossed components are required already at $(0,3/4,3/4)$.

\section{Conclusion}

The correct primary object is the free-field matrix element
\eqref{eq:direct-matrix-element}.  For the first crossed branch its
two-screening evaluation gives the irreducible $H$ restriction
\eqref{eq:N2-H}, and its natural three-screening evaluation gives the
factorized $K$ restriction \eqref{eq:N3-K}.  These are genuine
Pfaffian--Selberg calculations and show directly that the factor left after
the four-$\ell$ term is removed can be momentum dependent.

The finite function \eqref{eq:literal-A} remains useful as a PBW/Ward oracle
and has passed the stated low-level comparisons, but it is not a derivation
of the contour integral.

The status of the residue method is more precise.  In an ordinary scalar
Selberg family, endpoint support fixes the Gamma divisor and a collision
residue fixes the remaining normalization from the $N=0,1$ base cases.  The
Ramond descendant integral is vector-valued, and the two endpoint residues
proved here contain complementary rank-one projectors.  Neither relation
alone determines the full function.  A generic all-level answer requires a
proof of finite contour-orbit closure, degree and pole bounds on that space,
a full-rank collection of endpoint/cluster maps, and the inversion/crossing
connection matrix.  These data must include the current--fermion terms of
the reflected Ramond vertex and have no common invisible subspace.  Once
the vector Selberg integral is fixed from its base vectors, enough
same-parity screening nodes can separately reconstruct the pole-cleared
generic-momentum function by finite interpolation.

Consequently this note proposes the multi-channel screening organization
and proves its first nontrivial restrictions; it does not claim a completed
generic-momentum Ramond branching formula.  The projectors are auxiliary
data in an evaluation algorithm, not additional physical branching
coefficients.

\begin{thebibliography}{9}
\bibitem{HJ}
L.~Hadasz and Z.~Jaskolski,
\emph{Super-Liouville -- double Liouville correspondence},
JHEP 05 (2014) 124,
\href{https://arxiv.org/abs/1312.4520}{arXiv:1312.4520}.

\bibitem{BS}
M.~Bershtein and A.~Shchechkin,
\emph{B\"acklund transformation of Painlev\'e III$(D_8)$ tau function},
J.~Phys.~A 50 (2017) 115205,
\href{https://arxiv.org/abs/1608.02568}{arXiv:1608.02568}.

\bibitem{SS}
V.~Schomerus and P.~Suchanek,
\emph{Liouville's imaginary shadow},
JHEP 12 (2012) 020,
\href{https://arxiv.org/abs/1210.1856}{arXiv:1210.1856}.
\end{thebibliography}

\end{document}
